\section{Sujet TI n\textsuperscript{o}\,1 — Bibliothèque scolaire}

\begin{synthese}[Un sujet type d'examen]
Voici un \textbf{sujet type} de Baccalauréat, série \textbf{TI}, épreuve
\textbf{Algorithmique et Programmation}, calqué sur la structure réelle de
l'examen. Il réunit une \textbf{Partie I — Programmation web} (HTML, JavaScript,
PHP, MySQL) et une \textbf{Partie II — Programmation procédurale en C}
(structures, fonctions, trace). Contexte : la gestion informatisée de la
\textbf{bibliothèque d'un lycée camerounais}. Compose d'abord seul, dans les
conditions de l'examen, puis confronte ta copie au corrigé détaillé qui suit.
\end{synthese}

% =====================================================================
%  EN-TÊTE DU SUJET
% =====================================================================
\begin{center}
\fbox{\begin{minipage}{0.93\linewidth}
\centering\sffamily
{\bfseries OFFICE DU BACCALAURÉAT DU CAMEROUN}\\[1pt]
{\footnotesize Examen : Baccalauréat de l'Enseignement Secondaire Technique}\\[3pt]
\rule{0.6\linewidth}{0.4pt}\\[3pt]
{\large\bfseries SUJET TYPE — Session de juin}\\[2pt]
{\bfseries Série : TI}\\[3pt]
{\bfseries Épreuve : Algorithmique et Programmation}\\[2pt]
\begin{tabular}{@{}l@{\quad}c@{\quad}l@{}}
Durée : \textbf{03 heures} & $\bullet$ & Coefficient : \textbf{04}
\end{tabular}\\[3pt]
\rule{0.6\linewidth}{0.4pt}\\[2pt]
{\footnotesize\itshape L'épreuve comporte deux parties indépendantes, notées
chacune sur 10 points. L'usage de la calculatrice n'est pas autorisé.}
\end{minipage}}
\end{center}

\medskip
\noindent\textbf{Contexte général.} \emph{« Le Lycée Bilingue de Bafoussam
souhaite informatiser la gestion de sa bibliothèque scolaire, jusque-là tenue
sur des registres papier. L'établissement confie à un stagiaire en informatique
le développement d'une petite application web de gestion du fonds documentaire :
une page permet de saisir les caractéristiques d'un livre, et ces informations
sont enregistrées dans une base de données nommée \code{BIBLIO\_DB}, dont la
table principale est nommée \code{LIVRES} (champs : titre, auteur, catégorie,
nombre d'exemplaires). En parallèle, un programme en langage C est écrit pour
gérer hors ligne la liste des ouvrages. »}

% =====================================================================
%  PARTIE I — PROGRAMMATION WEB
% =====================================================================
\subsection*{Partie I — Programmation web \hfill (10 points)}
\addcontentsline{toc}{subsection}{Sujet TI nº1 — Partie I : Programmation web}

\subsubsection*{Exercice 1 — Programmation HTML et JavaScript \hfill (03,5 pts)}

\begin{enumerate}[label=\textbf{\arabic*.}]
\item \textbf{(0,25 pt)} Définir : \emph{page web statique}.
\item \textbf{(0,25 pt)} Citer le langage qui permet de structurer le contenu
d'une page web et celui qui permet de la rendre interactive côté client.
\item \textbf{(1,5 pt)} Écrire le code HTML affichant le formulaire de saisie
d'un livre (nommé \code{bookForm}). La méthode est \code{POST} et le fichier de
traitement est nommé \code{ajouter\_livre.php}. Les champs sont respectivement
\emph{« Titre »}, \emph{« Auteur »}, \emph{« Catégorie »}, \emph{« Nombre
d'exemplaires »} et le bouton \emph{« Ajouter »}. On utilisera un tableau de 05
lignes et 02 colonnes. Le formulaire appellera, au clic du bouton, une fonction
JavaScript de validation nommée \code{validateForm()}.
\item \textbf{(1,5 pt)} On veut, au clic du bouton \emph{« Ajouter »}, vérifier
que les champs \emph{Titre}, \emph{Auteur} et \emph{Catégorie} ne sont pas vides,
et que le champ \emph{Nombre d'exemplaires} contient bien une valeur numérique
(et non du texte). Si tout est correct, une boîte de dialogue affiche
\emph{« Livre prêt à être enregistré »} ; sinon, elle affiche
\emph{« Saisie invalide : vérifiez les champs »} et l'envoi du formulaire est
annulé. Écrire en JavaScript le code de la fonction \code{validateForm()}.
\end{enumerate}

\subsubsection*{Exercice 2 — Programmation PHP et MySQL \hfill (06,5 pts)}

\noindent On s'intéresse maintenant à la base de données \code{BIBLIO\_DB} et au
moyen d'y accéder depuis le script de traitement \code{ajouter\_livre.php}.

\begin{enumerate}[label=\textbf{\arabic*.}]
\item \textbf{(0,25 pt)} Définir : \emph{base de données}.
\item \textbf{(0,25 pt)} Donner la signification du sigle \emph{SGBD}.
\item \textbf{(0,25 pt)} Indiquer le côté (client ou serveur) où s'exécute un
script PHP.
\item \textbf{(0,75 pt)} Donner les éléments constitutifs d'un environnement
intégré de développement web dynamique (citer au moins trois éléments).
\item \textbf{(1 pt)} Donner le rôle des fonctions suivantes :
\code{mysqli\_connect()} ; \code{mysqli\_query()}.
\item \textbf{(0,5 pt)} Présenter la différence entre les fonctions
\code{mysqli\_fetch\_assoc()} et \code{mysqli\_fetch\_row()}.
\item \textbf{(1,75 pt)} Écrire le code PHP du fichier \code{ajouter\_livre.php}
qui : (a) se connecte à la base \code{BIBLIO\_DB} (hôte \code{localhost},
utilisateur \code{root}, mot de passe vide) ; (b) récupère les données envoyées
par le formulaire ; (c) les affiche dans un tableau HTML (une ligne par champ).
\item \textbf{(1,75 pt)} Écrire la requête SQL, puis le code PHP, permettant
\textbf{d'insérer} le nouveau livre dans la table \code{LIVRES} à partir des
données récupérées du formulaire, en signalant le succès ou l'échec de
l'opération.
\end{enumerate}

% =====================================================================
%  PARTIE II — PROGRAMMATION PROCÉDURALE EN C
% =====================================================================
\subsection*{Partie II — Programmation procédurale en C \hfill (10 points)}
\addcontentsline{toc}{subsection}{Sujet TI nº1 — Partie II : Programmation en C}

\subsubsection*{Exercice 1 \hfill (5 pts)}

\noindent\emph{« Pour gérer la bibliothèque hors ligne, on caractérise chaque
livre par un code (entier), un titre (chaîne), un auteur (chaîne) et un nombre
d'exemplaires (entier). On souhaite ranger l'ensemble des ouvrages dans un
tableau, puis disposer d'une fonction permettant de retrouver rapidement un
livre à partir de son code. »} Réaliser les tâches ci-dessous.

\begin{enumerate}[label=\textbf{\arabic*.}]
\item \textbf{(1,5 pt)} Donner l'instruction C permettant de déclarer une
structure d'enregistrement nommée \code{Livre} contenant les caractéristiques
décrites dans le texte (les chaînes seront des tableaux de 50 caractères).
\item \textbf{(0,5 pt)} Donner l'instruction C permettant de déclarer un tableau
\code{biblio} de 200 éléments de type \code{struct Livre}.
\item \textbf{(0,5 pt)} Donner deux avantages de l'utilisation des fonctions
dans un programme.
\item \textbf{(2,5 pts)} Donner en C le code de définition de la fonction
\code{rechercherLivre} qui recherche un livre par son code dans le tableau, en
utilisant la \textbf{recherche séquentielle}. Si le livre est trouvé, la
fonction affiche son titre, son auteur et son nombre d'exemplaires, puis renvoie
l'indice de la case trouvée ; sinon elle affiche le message \emph{« Aucun livre
ne correspond à ce code »} et renvoie $-1$. L'en-tête imposé de la fonction est :
\code{int rechercherLivre(struct Livre biblio[], int n, int code)}.
\end{enumerate}

\subsubsection*{Exercice 2 \hfill (5 pts)}

\noindent On donne le programme C ci-dessous, qui calcule la \textbf{somme des
chiffres} d'un entier naturel saisi par l'utilisateur (les numéros de ligne 1 à
22 sont rappelés en commentaire pour pouvoir y faire référence) :

\begin{codec}
/* 1*/  /* Calcul de la somme des chiffres d'un entier naturel */
/* 2*/  #include <stdio.h>
/* 3*/  #include <stdlib.h>
/* 4*/  int sommeChiffres(int n);
/* 5*/  int main() {
/* 6*/      int x, s;
/* 7*/      printf("Entrez un entier naturel : ");
/* 8*/      scanf("%d", &x);
/* 9*/      if (x < 0) {
/*10*/          printf("Nombre negatif non autorise !\n");
/*11*/      }
/*12*/      else {
/*13*/          s = sommeChiffres(x);
/*14*/          printf("Somme des chiffres de %d = %d\n", x, s);
/*15*/      }
/*16*/      return 0;
/*17*/  }
/*18*/  int sommeChiffres(int n) {
/*19*/      int somme = 0;
/*20*/      while (n != 0) { somme = somme + n % 10; n = n / 10; }
/*21*/      return somme;
/*22*/  }
\end{codec}

\begin{enumerate}[label=\textbf{\arabic*.}]
\item Donner le rôle de :
  \begin{enumerate}[label=\textbf{1.\arabic*}]
  \item \textbf{(0,25 pt)} la ligne 4 (la déclaration
  \code{int sommeChiffres(int n);}) ;
  \item \textbf{(0,25 pt)} la ligne 8 (\code{scanf("\%d", \&x);}).
  \end{enumerate}
\item \textbf{(0,5 pt)} Présenter la différence fondamentale entre la variable
\code{x} (ligne 6) et la variable \code{n} (ligne 18) dans ce programme.
\item \textbf{(1,5 pt)} Exécuter pas à pas la fonction \code{sommeChiffres}
lorsque \code{x = 472} (présenter la trace dans un tableau : tours de boucle,
valeurs de \code{n}, de \code{n \% 10} et de \code{somme}).
\item \textbf{(1 pt)} Réécrire les instructions de la structure \code{if/else}
(lignes 9 à 15) en utilisant l'opérateur ternaire (\code{ ? : }) pour déterminer
le message, l'affichage de la somme restant inchangé.
\item \textbf{(1,5 pt)} Réécrire la fonction \code{sommeChiffres} sous forme
\textbf{récursive}.
\end{enumerate}

% =====================================================================
%  CORRIGÉ
% =====================================================================
\corrige{du sujet TI n\textsuperscript{o}\,1}

\subsubsection*{Partie I — Exercice 1 (HTML / JavaScript)}

\textbf{1. Définition — page web statique.} Une \emph{page web statique} est une
page dont le contenu est figé : il est écrit une fois pour toutes en HTML et
s'affiche à l'identique pour tous les visiteurs, sans traitement côté serveur ni
accès à une base de données. (Par opposition, une page \emph{dynamique} est
générée à la volée, par exemple par PHP.)

\medskip
\textbf{2. Langages.} Le langage qui structure le contenu d'une page web est
\textbf{HTML} ; celui qui la rend interactive côté client est \textbf{JavaScript}.

\medskip
\textbf{3. Formulaire HTML (\code{bookForm}) avec tableau 5$\times$2.}

\begin{codehtml}
<!DOCTYPE html>
<html lang="fr">
<head>
  <meta charset="utf-8">
  <title>Ajout d'un livre</title>
  <script src="validation.js"></script>
</head>
<body>
  <h2>Nouveau livre</h2>
  <form name="bookForm" method="POST" action="ajouter_livre.php"
        onsubmit="return validateForm()">
    <table border="1" cellpadding="6">
      <tr>
        <td>Titre</td>
        <td><input type="text" name="titre" id="titre"></td>
      </tr>
      <tr>
        <td>Auteur</td>
        <td><input type="text" name="auteur" id="auteur"></td>
      </tr>
      <tr>
        <td>Catégorie</td>
        <td><input type="text" name="categorie" id="categorie"></td>
      </tr>
      <tr>
        <td>Nombre d'exemplaires</td>
        <td><input type="text" name="nbExemplaires" id="nbExemplaires"></td>
      </tr>
      <tr>
        <td colspan="2" align="center">
          <input type="submit" name="ajouter" value="Ajouter">
        </td>
      </tr>
    </table>
  </form>
</body>
</html>
\end{codehtml}

\fig{Tableau de 5 lignes $\times$ 2 colonnes ; la validation est déclenchée par \code{onsubmit}.}

\begin{remarque}
On donne à chaque champ un attribut \code{id} pour pouvoir le retrouver facilement
en JavaScript avec \code{document.getElementById}. Le champ
\emph{nombre d'exemplaires} est de type \code{text} (et non \code{number}) afin de
pouvoir contrôler nous-mêmes, en JavaScript, qu'une valeur numérique a bien été
saisie.
\end{remarque}

\medskip
\textbf{4. Fonction \code{validateForm()}.}

\begin{codejs}
function validateForm() {
  // Récupération des valeurs saisies (sans espaces superflus)
  var titre  = document.getElementById("titre").value.trim();
  var auteur = document.getElementById("auteur").value.trim();
  var categorie = document.getElementById("categorie").value.trim();
  var nbExp  = document.getElementById("nbExemplaires").value.trim();

  // 1) Les trois champs texte doivent etre non vides
  if (titre === "" || auteur === "" || categorie === "") {
    alert("Saisie invalide : verifiez les champs");
    return false;   // annule l'envoi du formulaire
  }

  // 2) Le nombre d'exemplaires doit etre une valeur numerique
  if (nbExp === "" || isNaN(nbExp)) {
    alert("Saisie invalide : verifiez les champs");
    return false;
  }

  // Tout est correct
  alert("Livre pret a etre enregistre");
  return true;    // autorise l'envoi vers ajouter_livre.php
}
\end{codejs}

\begin{remarque}
\code{isNaN(nbExp)} renvoie \code{true} quand la valeur \emph{n'est pas un
nombre}. En renvoyant \code{false}, la fonction empêche l'envoi du formulaire
(grâce à \code{onsubmit="return validateForm()"}).
\end{remarque}

\subsubsection*{Partie I — Exercice 2 (PHP / MySQL)}

\textbf{1. Définition — base de données.} Une \emph{base de données} est un
ensemble structuré et organisé de données, stockées de façon durable, de manière
à pouvoir être facilement recherchées, consultées et mises à jour.

\medskip
\textbf{2. SGBD.} \emph{Système de Gestion de Base de Données} (en anglais
\emph{DBMS}). Exemples : MySQL, PostgreSQL, Oracle.

\medskip
\textbf{3. Côté d'exécution de PHP.} Un script PHP s'exécute du \textbf{côté
serveur} ; le navigateur du client ne reçoit que le résultat (du HTML).

\medskip
\textbf{4. Éléments d'un environnement intégré de développement web.} On retrouve
au minimum : un \textbf{serveur web} (Apache), un \textbf{interpréteur PHP} et un
\textbf{serveur de base de données} (MySQL/MariaDB). Une suite comme
\emph{WampServer} ou \emph{XAMPP} regroupe ces trois éléments ; on y ajoute en
général un éditeur de code et l'outil d'administration \emph{phpMyAdmin}.

\medskip
\textbf{5. Rôle des fonctions.}
\begin{itemize}
\item \code{mysqli\_connect()} : ouvre une connexion au serveur MySQL et
sélectionne la base de données ; elle renvoie un identifiant de connexion.
\item \code{mysqli\_query()} : envoie (exécute) une requête SQL sur la base via
la connexion ouverte ; pour un \code{SELECT}, elle renvoie un jeu de résultats.
\end{itemize}

\medskip
\textbf{6. Différence \code{mysqli\_fetch\_assoc()} / \code{mysqli\_fetch\_row()}.}
Les deux extraient une ligne du résultat sous forme de tableau, mais :
\code{mysqli\_fetch\_assoc()} renvoie un tableau \textbf{associatif} indexé par
le \emph{nom des colonnes} (ex. \code{\$row["titre"]}), tandis que
\code{mysqli\_fetch\_row()} renvoie un tableau \textbf{indexé numériquement} à
partir de 0 (ex. \code{\$row[0]}).

\medskip
\textbf{7. \code{ajouter\_livre.php} — connexion, récupération, affichage.}

\begin{codephp}
<?php
  // (a) Connexion a la base BIBLIO_DB
  $conn = mysqli_connect("localhost", "root", "", "BIBLIO_DB");
  if (!$conn) {
    die("Echec de la connexion : " . mysqli_connect_error());
  }

  // (b) Recuperation des donnees envoyees par le formulaire (methode POST)
  $titre        = $_POST["titre"];
  $auteur       = $_POST["auteur"];
  $categorie    = $_POST["categorie"];
  $nbExemplaires = $_POST["nbExemplaires"];

  // (c) Affichage des donnees dans un tableau HTML
  echo "<table border='1' cellpadding='6'>";
  echo "<tr><td>Titre</td><td>" . $titre . "</td></tr>";
  echo "<tr><td>Auteur</td><td>" . $auteur . "</td></tr>";
  echo "<tr><td>Categorie</td><td>" . $categorie . "</td></tr>";
  echo "<tr><td>Nombre d'exemplaires</td><td>" . $nbExemplaires . "</td></tr>";
  echo "</table>";
?>
\end{codephp}

\medskip
\textbf{8. Requête SQL d'insertion + code PHP.} La requête SQL d'insertion est :

\begin{codesql}
INSERT INTO LIVRES (titre, auteur, categorie, nbExemplaires)
VALUES ('Things Fall Apart', 'Chinua Achebe', 'Roman', 12);
\end{codesql}

\noindent En PHP, en réutilisant les données récupérées du formulaire :

\begin{codephp}
<?php
  // Construction de la requete avec les valeurs du formulaire
  $sql = "INSERT INTO LIVRES (titre, auteur, categorie, nbExemplaires)
          VALUES ('$titre', '$auteur', '$categorie', '$nbExemplaires')";

  // Execution de la requete
  if (mysqli_query($conn, $sql)) {
    echo "<p>Livre ajoute avec succes !</p>";
  } else {
    echo "<p>Echec de l'ajout : " . mysqli_error($conn) . "</p>";
  }

  mysqli_close($conn);   // fermeture de la connexion
?>
\end{codephp}

\begin{remarque}
Dans une application réelle, on protégerait la requête contre les injections SQL
(requêtes préparées, \code{mysqli\_real\_escape\_string}). Au Bac TI, l'écriture
directe ci-dessus est acceptée.
\end{remarque}

\subsubsection*{Partie II — Exercice 1 (structures et fonctions en C)}

\textbf{1. Déclaration de la structure \code{Livre}.}

\begin{codec}
struct Livre {
    int  code;             /* code du livre (entier)        */
    char titre[50];        /* titre (chaine de 50 car.)     */
    char auteur[50];       /* auteur (chaine de 50 car.)    */
    int  nbExemplaires;    /* nombre d'exemplaires (entier) */
};
\end{codec}

\medskip
\textbf{2. Déclaration du tableau \code{biblio}.}

\begin{codec}
struct Livre biblio[200];
\end{codec}

\medskip
\textbf{3. Deux avantages des fonctions.}
\begin{itemize}
\item Elles évitent la répétition du code : on écrit un traitement une seule fois
et on l'appelle autant de fois que nécessaire (réutilisabilité).
\item Elles rendent le programme plus lisible et plus facile à corriger, en le
découpant en blocs courts ayant chacun un rôle précis (modularité).
\end{itemize}

\medskip
\textbf{4. Fonction \code{rechercherLivre} (recherche séquentielle).}

\begin{codec}
int rechercherLivre(struct Livre biblio[], int n, int code) {
    int i;
    for (i = 0; i < n; i++) {
        if (biblio[i].code == code) {        /* livre trouve */
            printf("Titre   : %s\n", biblio[i].titre);
            printf("Auteur  : %s\n", biblio[i].auteur);
            printf("Exempl. : %d\n", biblio[i].nbExemplaires);
            return i;                        /* indice de la case trouvee */
        }
    }
    printf("Aucun livre ne correspond a ce code\n");
    return -1;                               /* non trouve */
}
\end{codec}

\begin{remarque}
On parcourt le tableau case par case (de l'indice 0 à \code{n-1}) : c'est la
\emph{recherche séquentielle}. Dès que le code recherché est rencontré, on
affiche les informations et on quitte la fonction avec \code{return}, sans
parcourir le reste du tableau.
\end{remarque}

\subsubsection*{Partie II — Exercice 2 (programme « somme des chiffres »)}

\textbf{1. Rôle des lignes.}
\begin{itemize}
\item \textbf{Ligne 4} (\code{int sommeChiffres(int n);}) : c'est le
\emph{prototype} (déclaration) de la fonction. Il annonce au compilateur le nom
de la fonction, le type de son paramètre (\code{int}) et le type de sa valeur de
retour (\code{int}), avant son utilisation dans \code{main}.
\item \textbf{Ligne 8} (\code{scanf("\%d", \&x);}) : elle lit au clavier un
entier saisi par l'utilisateur et le range dans la variable \code{x} (le
\code{\&} fournit l'adresse de \code{x}).
\end{itemize}

\medskip
\textbf{2. Différence entre \code{x} et \code{n}.} \code{x} est une variable
\textbf{locale à \code{main}} : elle contient la valeur saisie par l'utilisateur.
\code{n} est le \textbf{paramètre formel} de la fonction \code{sommeChiffres} :
c'est une variable locale à la fonction qui reçoit une \emph{copie} de la valeur
de \code{x} au moment de l'appel (passage par valeur). Modifier \code{n} dans la
fonction ne change donc pas \code{x} dans \code{main}.

\medskip
\textbf{3. Trace pour \code{x = 472}.} La fonction est appelée avec \code{n = 472}
et \code{somme = 0}. La boucle \code{while} tourne tant que \code{n != 0} :

\begin{center}\small
\begin{tabular}{@{}ccccc@{}}
\toprule
\textbf{Tour} & \textbf{n (début)} & \textbf{n \% 10} & \textbf{somme} & \textbf{n (fin = n / 10)}\\
\midrule
1 & 472 & 2 & 0 + 2 = 2  & 47\\
2 & 47  & 7 & 2 + 7 = 9  & 4\\
3 & 4   & 4 & 9 + 4 = 13 & 0\\
\bottomrule
\end{tabular}
\end{center}

\noindent À l'issue du tour 3, \code{n} vaut 0 : la boucle s'arrête. La fonction
renvoie \code{somme = 13}. Le programme affiche :
\emph{« Somme des chiffres de 472 = 13 »}.

\medskip
\textbf{4. Réécriture du \code{if/else} avec l'opérateur ternaire.} Le bloc des
lignes 9 à 15 peut se réécrire ainsi (on choisit le message avec le ternaire) :

\begin{codec}
char *message = (x < 0) ? "Nombre negatif non autorise !"
                        : "OK";
if (x < 0)
    printf("%s\n", message);
else {
    s = sommeChiffres(x);
    printf("Somme des chiffres de %d = %d\n", x, s);
}
\end{codec}

\noindent Une forme plus compacte, lorsqu'on ne veut afficher qu'un message :

\begin{codec}
(x < 0) ? printf("Nombre negatif non autorise !\n")
        : printf("Somme des chiffres de %d = %d\n", x, sommeChiffres(x));
\end{codec}

\medskip
\textbf{5. Version récursive de \code{sommeChiffres}.}

\begin{codec}
int sommeChiffres(int n) {
    if (n == 0)                       /* cas de base */
        return 0;
    return (n % 10) + sommeChiffres(n / 10);   /* dernier chiffre + reste */
}
\end{codec}

\begin{remarque}
À chaque appel, on isole le dernier chiffre avec \code{n \% 10}, puis on relance
la fonction sur \code{n / 10} (le nombre privé de son dernier chiffre).
La récursion s'arrête quand \code{n} vaut 0 : c'est le \emph{cas de base},
indispensable pour éviter une récursion infinie.
\end{remarque}

\begin{aretenir}
Un sujet TI se compose toujours de deux parties indépendantes notées sur 10.
\textbf{Partie web} : sache écrire un formulaire HTML dans un tableau, une
fonction \code{validateForm()} (champs non vides, \code{isNaN}, \code{return
false}), une connexion \code{mysqli\_connect} et une requête (\code{INSERT},
\code{SELECT}, \code{DELETE}). \textbf{Partie C} : maîtrise la déclaration d'une
\code{struct}, un tableau de structures, la recherche séquentielle, la trace
pas-à-pas d'une boucle et le passage d'une version itérative à une version
récursive.
\end{aretenir}
